\section{Introducción}

% ESQUELETO (Eva, 06/11). Las líneas con % son el guion de cada apartado:
% no salen en el PDF. En la semana del 9-15/11 se convierten en texto.
% Ningún número de la red se escribe aquí a mano: salen de results/.

\subsection{La disección aórtica}
% Definición: desgarro de la íntima con separación de las capas de la media
%   y formación de una falsa luz \cite{Nienaber2016}.
% Clasificación de Stanford (A/B) y de DeBakey \cite{Isselbacher2022}.
% Epidemiología: incidencia poblacional y mortalidad antes de llegar
%   al hospital \cite{Howard2013}.
% Evolución clínica y mortalidad hospitalaria (registro IRAD) \cite{Evangelista2018}.
% Factor de riesgo principal: hipertensión \cite{Howard2013}.

\subsection{Fisiopatología de la pared aórtica}
% Degeneración de la media: matriz extracelular (fibrilina, colágeno, elastina).
% Células de músculo liso y su aparato contráctil.
% Señalización de TGF-$\beta$ \cite{Pinard2019}.

\subsection{Base genética}
%  Formas sindrómicas: Marfan (FBN1), Loeys-Dietz (TGFBR1/2, SMAD3, TGFB2/3),
%   Ehlers-Danlos vascular (COL3A1) \cite{Pinard2019}.
% Formas no sindrómicas: ACTA2, MYH11, MYLK, PRKG1 \cite{Renard2018}.
% Clasificación de genes por evidencia clínica (ClinGen) \cite{Renard2018}.

\subsection{El término HPO HP:0002647 (\textit{Aortic dissection})}
% Qué es la Human Phenotype Ontology y cómo asocia fenotipos y genes \cite{Gargano2024}.
% - Por qué partimos de un fenotipo y no de una enfermedad concreta.
% - NO poner aquí el número de genes: va en Resultados (Aitana), sacado de results/.

\subsection{Biología de sistemas y medicina de redes}
% Una enfermedad como perturbación de módulos en una red, no de un único gen \cite{Barabasi2011}.
% Redes de interacción proteína-proteína: STRING \cite{Szklarczyk2023}.

\subsection{Objetivos}
% 1. Obtener los genes asociados a HP:0002647.
% 2. Construir y analizar su red de interacción (topología).
% 3. Identificar módulos y sus funciones (enriquecimiento).
% (Revisar cuando Antonio y Nico cierren el análisis.)
